% ============================================================================
% PRIMEFIELD: THE THERMODYNAMIC TRIE
% ============================================================================
\subsection{Contiguous GPU-Dispatch Array}
\label{sec:primefield_layout}

The PrimeField is the compute-side representation of memory: a flat,
contiguous array of \texttt{IcosahedralManifold} structures that the GPU
scans in parallel.  Each manifold occupies a fixed $8.64$~KB regardless
of fill state, guaranteeing stable pointer arithmetic and aligned warp
strides.

\begin{itemize}[nosep]
  \item \textbf{Insertion:}  Tokens are merged into the active manifold
        via \texttt{ChordOR} at a block index determined by the token's
        topological events and chord dynamics.
  \item \textbf{Rotation:}  Structural events physically permute the
        active manifold's block array.  Because $A_5$ permutations are
        applied as array index remappings, they execute in $O(1)$ without
        data copying.
  \item \textbf{Atomic Snapshots:}  Read access returns a \texttt{(pointer,\,N)}
        pair under atomic load, preventing concurrency tearing during
        background GPU dispatch.
\end{itemize}

\subsection{Morton-Coded Spatial Indexing}
\label{sec:morton_codes}

Each token is assigned a deterministic 64-bit \emph{Morton key} that packs
three coordinates into a single integer for locality-preserving sort order:

\begin{equation}
  \mathcal{M}(z,\,\text{sym},\,\text{pos})
    = \bigl(\texttt{uint64}(z) \ll 56\bigr)
    \;\big|\;
      \bigl(\texttt{uint64}(\text{sym}) \ll 32\bigr)
    \;\big|\;
      \texttt{uint64}(\text{pos})
  \label{eq:morton_key}
\end{equation}

\noindent The bit-field layout from MSB to LSB is:

\begin{center}
\begin{tabular}{ccc}
  \toprule
  \textbf{Bits 63--56} & \textbf{Bits 55--32} & \textbf{Bits 31--0} \\
  \midrule
  $z$ (Scale Index) & Symbol (Byte Value) & Position (Sequence Offset) \\
  8 bits & 24 bits & 32 bits \\
  \bottomrule
\end{tabular}
\end{center}

This layout is load-bearing in three ways:
\begin{enumerate}[nosep]
  \item \textbf{Scale stratification:}  Grouping by $z$ in the MSBs
        ensures that all entries at the same hierarchical depth sort
        contiguously, enabling $O(\log N)$ binary search within a
        scale layer.
  \item \textbf{Symbol locality:}  Within a scale, all occurrences of the
        same byte value occupy a contiguous key range
        $[\text{sym} \ll 32,\; (\text{sym}+1) \ll 32 - 1]$, allowing
        \texttt{QueryByByte} to narrow a 50M-entry store to a single
        binary-search band.
  \item \textbf{Positional contiguity:}  Within a symbol band, positions
        are numerically ordered in the LSBs, making neighborhood
        queries (\texttt{QueryNeighborhood}) a simple range scan.
\end{enumerate}

\subsection{Zipfian Distributions and Collision-as-Compression}
\label{sec:collision_is_compression}

Natural data sources---text, images, audio---universally exhibit
\emph{Zipfian} frequency distributions: a small number of byte-position
patterns account for the overwhelming majority of occurrences.  The LSM
Spatial Index exploits this directly.

When two tokens share the same Morton key (identical byte value at the
same scale and position), the LSM merge operation natively \emph{coalesces}
them into a single entry.  This is not a hash collision requiring
resolution---it is lossless deduplication:

\begin{equation}
  \text{Compression Ratio} =
    \frac{\text{Raw Token Count}}{\text{Unique Morton Keys}}
    \;=\;
    \frac{N}{\lvert\{\mathcal{M}(z_i, s_i, p_i)\}_{i=1}^{N}\rvert}
  \label{eq:compression_ratio}
\end{equation}

Because Zipf's Law guarantees that the frequency of the $r$-th most common
pattern scales as $f(r) \propto r^{-\alpha}$ with $\alpha \approx 1$,
the number of unique keys grows sub-linearly with corpus size.
Empirically, we observe that compression ratios \textbf{increase} as
datasets grow larger, because larger corpora sample the tail of the
Zipfian distribution more densely, producing proportionally more
collisions per novel entry.

\begin{table}[h]
  \centering
  \caption{Collision-compression metrics across datasets.  Ratios increase
           with dataset size, consistent with Zipfian redundancy scaling.}
  \label{tab:compression_metrics}
  \begin{tabular}{lrrrr}
    \toprule
    \textbf{Dataset}
      & \textbf{Raw Tokens}
      & \textbf{Unique Keys}
      & \textbf{Dedup.\ \%}
      & \textbf{Compression $\times$} \\
    \midrule
    CIFAR-10 (60K images)
      & 195,000,000
      & 3,693,182
      & 98.1\%
      & $52.8\times$ \\
    \addlinespace
    \textit{MBPP (Python)}
      & \textit{---}
      & \textit{---}
      & \textit{---}
      & \textit{---} \\
    \textit{bAbI (QA)}
      & \textit{---}
      & \textit{---}
      & \textit{---}
      & \textit{---} \\
    \textit{Combined Corpus}
      & \textit{---}
      & \textit{---}
      & \textit{---}
      & \textit{---} \\
    \bottomrule
  \end{tabular}
\end{table}

\noindent The theoretical lower bound on compression is the dataset's
Shannon entropy $H$: the system cannot compress below the information
content of the source distribution.  The gap between observed compression
and the entropy bound measures the efficiency of the Morton key design.

\subsection{Entropy-Driven Sequence Boundaries}
\label{sec:sequence_boundaries}

When the Sequencer detects a topological boundary and the active manifold
carries signal (non-empty chords), the PrimeField freezes the current
manifold and advances to a new one.  The Z-coordinate of the Morton key
tracks scale depth: each sequence reset increments $z$, stratifying memory
into episodic layers.

\subsection{Prototype Snapping and Cleanup}
\label{sec:prototype_snapping}

To stabilize the memory substrate against noise, the PrimeField maintains
per-block prototype attractors.  Incoming chords are compared against stored
prototypes; if similarity exceeds a threshold, the chord is ``snapped'' to
the prototype---reinforcing established patterns.  Novel chords that diverge
sufficiently seed new prototypes, enabling the substrate to discover and
crystallize recurring structural motifs.

\subsection{Manifold Masking for Self-Exclusion}
\label{sec:masking}

During BestFill search, the system can \texttt{Mask} individual manifolds
(temporarily zeroing them) to exclude recently matched spans from
consideration.  This prevents the popcount echo-attractor problem where the
same high-scoring span is retrieved repeatedly.  Masked manifolds are restored
via \texttt{Unmask} after the search completes.

